
\AddSymbolsF
\pagecolor{ForestGreen!15!Gray!}
\color{white}
\quad\quad \lettrine[,lraise=0.2,lhang=0.5]
{\textbf{伊}}{藤}佐一听完教授的话，仿佛若有所思——回味着，又观察着，紧接着便代入了自己。他又想起自己这粒灰尘；也回想起远方那粒灰尘——可是那灰尘，在他心目中，分量是有多大啊。

\twkkai{我们去别的地方逛逛？}教授忽然用手撑着草地，直起身子，又拍了拍大腿外侧——准备站起。

\twkkai{哦——好的；}伊藤佐一应和着。

他们都站了起来，看向那仍纠缠在一起的风筝。教授低头又擦了擦眼镜，接着突然拍了拍伊藤佐一的肩膀。

\twkkai{你的东西掉了；}他和蔼地说。

伊藤佐一看向草地——是他的奖章。他仿佛像泄露了什么似的，连忙弯下身捡起。他才再次发现上面的铭文。

\twkkai{对了，教授——如果你不介意……可以帮我看看吗？}

他又忘了他的担忧，大大方方地把奖章直接递给教授。

\twkkai{嗯哼……\(\forall \alpha \in \mathbb{R} \setminus \mathbb{Q},\ \exists^\infty \frac{p}{q} \in \mathbb{Q},\ \left| \alpha - \frac{p}{q} \right| < \frac{1}{q^2}.\)}

教授读着，接着说道。

\twkkai{啊——这是狄利克雷逼近定理啊。哪里不理解呢？}

伊藤佐一右手扯着衣服，看着教授T恤上的标志，发着呆，低声说着。

\twkkai{其实……我哪里都不理解——我连这个命题讲什么都不知道。}

教授把奖章归还给伊藤佐一，又转向另一边。

\twkkai{边走边聊吧。来，走这边；}他招了招手。伊藤佐一跟着。

\twkkai{奖章上的内容简直是完美适配我刚才灰尘的比喻……不过说来也奇怪，为什么每次都那么巧呢？}

\twkkai{什么那么巧？}伊藤佐一扯着衣服的右手攥得更紧了。

\twkkai{就是每次我感叹完，你给的奖章上的内容就会完美契合我感叹的内容；}教授看向伊藤佐一。\twkkai{话说，你究竟有多少个这样的奖章？}

\twkkai{也就两三个吧；}伊藤佐一撒着谎。他知道很不好受——但他实在没有办法。

\twkkai{说回正题吧。这条命题讲的就是——任意无理数\(\alpha\)，总有无穷多个有理数\(\frac{p}{q}\)能任意接近它，但永远不会真正等于它。}

\twkkai{就这样吗？}伊藤佐一似乎有点惊讶。

\twkkai{是的，就是这么简单。你又想起了什么吗？}教授看着伊藤佐一。

\twkkai{没有，没有……；}伊藤佐一眼睛死死盯着路面，回想着自己刚才在便利店的那些思考。

\twkkai{……自己和她，只是两粒灰尘，因风吸引着，又因风排斥着，永不能接触，又永不得安宁……}

\twkkai{永不得安宁……}

伊藤佐一想到这，不禁打了个寒颤。教授看到了；但他没说话。伊藤佐一尴尬得只想抽自己身体一巴掌。

\twkkai{要我讲讲证明过程吗——这个几乎不需要用草稿纸；}教授轻轻跳了一下，尝试够到经过他头顶的一片树枝——没够着。

\twkkai{嗯；}伊藤佐一答应了。他急需什么来转移自己的注意力。

\twkkai{好。那我们首先要做的，就是取定正整数\(N\)。接着考虑\(N+1\)个数：
$$\{0\cdot\alpha\}, \{1\cdot\alpha\}, \{2\cdot\alpha\}, \ldots, \{N\cdot\alpha\}$$
（其中\(\{x\}\)表示\(x\)的小数部分），它们都落在\([0,1)\)这个区间中。}

至此，伊藤佐一觉得一切正常。教授则接着讲了下去。

\twkkai{
把\([0,1)\)分成\(N\)个长度为\(\frac{1}{N}\)的小区间——你可以把它们想象成\(N\)块砖：\([0,1/N), [1/N,2/N), ..., [(N-1)/N,1)\)。
}

\twkkai{砖？}

伊藤佐一在心里问着教授。他知道，他又想起了便利店那些珍珠白色的砖，接着又是灰尘，最后又是——她。

他恳求教授讲快点；教授无意中配合着他——讲得越来越快了。

\twkkai{
\(N+1\)个数被放进了\(N\)个区间，那么，鸽巢原理保证至少有两个数落在同一个区间里——设它们是\(\{iα\}\)和\(\{jα\}, 0≤i<j≤N\)。
那么：
$$|\{j\alpha\} - \{i\alpha\}| < \frac{1}{N}$$
也就是存在整数\(p\)使得：
$$|(j-i)\alpha - p| < \frac{1}{N}$$
设\(q = j-i（1≤q≤N）\)，则：
$$|q\alpha - p| < \frac{1}{N} \le \frac{1}{q}$$}

\twkkai{跟上了吗？}教授突然问道——是那么地突兀。

伊藤佐一当然跟上了。他跟不上的是，为什么这么巧妙的证明，竟这么巧妙地与自己和她的关系相吻合。

\twkkai{嗯。}他爽快地回应了教授，以便腾出时间继续跟上自己——而不是教授——的思路。

\twkkai{好——于是，不等式两边除以\(q\)：
$$\left|\alpha - \frac{p}{q}\right| < \frac{1}{qN} \le \frac{1}{q^2}$$

这就是所要证明的内容。再让\(N\)不断增大，可以得到无穷多组这样的\((p,q)\)。}

伊藤佐一听完，非但没有松一口气，反而十分沮丧——没成想，到最后，教授的讲解，非但没让他转移注意力放松，还让他陷入自己的思绪更深了。

但这次，惊异于证明的精巧，伊藤佐一终于又对狄利克雷这个人物产生了兴趣。他终于在这一刻，问起了教授。

\twkkai{所以，教授——狄利克雷是谁？}

教授愣了一下，脚步踉跄着。他扶了扶眼镜。

\twkkai{有机会，我会带你去见他的。}

\twkkai{带我去见他？}

\twkkai{是的；不过，我也不知道能不能履行他派给我的职责。你要不要叫上铃木同学一起？}

伊藤佐一听到教授这么说，便摸了摸自己的通讯机。他犹豫着；但不过两秒，他便给铃木和子发了消息。

\twkkai{叮！}

\twkkai{回复了？}

教授仿佛比伊藤佐一更急切；而伊藤佐一点了点头。

\twkkai{好；你先约她出来，我得先去茶舍一趟。叫她直接去茶舍也行。你可以不跟着我。}

伊藤佐一迷惑地看着教授的背影。不等他思考，教授已经急匆匆地走远了。

远方的树枝上，多了几只洁白的鸽子。

它们正探着头，四处张望着。

它们不敢发出声音。

伊藤佐一则独自走在去往茶舍的小路上。他预感到了什么。





